% GENERATED FILE. Edit canonical lessons, then run npm run generate:books.
% canonical-source-hash: fnv1a64:ae398258ec5edb5a
% canonical-lessons: ZH-C261-jingyan, ZH-C261-cengjing, ZH-C261-liuxue-abroad, ZH-C261-guowai, ZH-C261-tiyan

\chapter{Have You Ever?}
\label{ch:zh-have-you-ever}
\hlchaptermodality{261}

\begin{chapteropening}
\textbf{By the end of this chapter:} \emph{I can talk about experience, say I once did something, study abroad, go overseas and say whom I saw.}

Complete the last lesson of chapter 261: I can talk about experience, say I once did something, study abroad, go overseas and say whom I saw.
\end{chapteropening}

\section[jīngyàn]{jīngyàn — experience; know-how}

\label{lesson:ZH-C261-jingyan}



\begin{quote}
Before the new one: say the Mandarin for technique, then the Mandarin for to make progress.
\end{quote}

\subsection*{You'll want to know: jīngyàn}
\emph{jīngyàn} — \textquotedblleft{}experience; know-how\textquotedblright{}.

\begin{quote}
\emph{tā hěn yǒu jīngyàn} — He's very experienced.
\end{quote}

Its characters come later; for now it is a word for the ear.

\begin{cousinweb}
\emph{jīngyàn} is built from \emph{jīng} and \emph{yàn}, both inherited from Middle Chinese, \emph{keng} and \emph{ngjemH} (Guangyun, in Baxter's transcription), which glosses the first 'constant' and 'the warp of cloth', and the second 'to verify'. How the pair came to mean experience is not settled.
\end{cousinweb}

\begin{grammarlens}[title={What you've built}]
One more word for talking about the past.
\end{grammarlens}

\subsection*{Your turn}
\begin{itemize}
\raggedright
  \item \emph{Say it:} \emph{jīngyàn}
  \item \emph{Say it:} \emph{jīngyàn}, once more
  \item \emph{Say it:} \emph{tā hěn yǒu jīngyàn} — He's very experienced.
  \item \emph{Recall:} say the Mandarin for technique, then the Mandarin for to make progress, then say \emph{jīngyàn} again
\end{itemize}

\begin{tcolorbox}[breakable,colback=teal!4,colframe=teal!35!black,title={Before you move on}]
What does jīngyàn mean? (\textquotedblleft{}Experience; know-how\textquotedblright{}.) Say it once more.
\end{tcolorbox}
\section[céngjīng]{céngjīng — once (in the past); at one time}

\label{lesson:ZH-C261-cengjing}



\begin{quote}
Before the new one: say the Mandarin for to make progress, then the Mandarin for know-how.
\end{quote}

\subsection*{You'll want to know: céngjīng}
\emph{céngjīng} — \textquotedblleft{}once (in the past); at one time\textquotedblright{}.

\begin{quote}
\emph{wǒ céngjīng qù guo Běijīng} — I've been to Beijing once.
\end{quote}

Its characters come later; for now it is a word for the ear.

\begin{cousinweb}
\emph{céngjīng} is built from \emph{céng} and \emph{jīng}, both inherited from Middle Chinese, \emph{dzong} and \emph{keng} (Guangyun, in Baxter's transcription), which glosses the first with the second, 'to pass through'.
\end{cousinweb}

\begin{grammarlens}[title={What you've built}]
One more word for talking about the past.
\end{grammarlens}

\subsection*{Your turn}
\begin{itemize}
\raggedright
  \item \emph{Say it:} \emph{céngjīng}
  \item \emph{Say it:} \emph{céngjīng}, once more
  \item \emph{Say it:} \emph{wǒ céngjīng qù guo Běijīng} — I've been to Beijing once.
  \item \emph{Recall:} say the Mandarin for to make progress, then the Mandarin for know-how, then say \emph{céngjīng} again
\end{itemize}

\begin{tcolorbox}[breakable,colback=teal!4,colframe=teal!35!black,title={Before you move on}]
What does céngjīng mean? (\textquotedblleft{}Once (in the past); at one time\textquotedblright{}.) Say it once more.
\end{tcolorbox}
\section[liúxué]{liúxué — to study abroad}

\label{lesson:ZH-C261-liuxue-abroad}



\begin{quote}
Before the new one: say the Mandarin for know-how, then the Mandarin for at one time.
\end{quote}

\subsection*{You'll want to know: liúxué}
\emph{liúxué} — \textquotedblleft{}to study abroad\textquotedblright{}.

\begin{quote}
\emph{wǒ xiǎng qù Zhōngguó liúxué} — I'd like to go to China to study.
\end{quote}

Its characters come later; for now it is a word for the ear.

\begin{cousinweb}
\emph{liúxué} is built from \emph{liú} 'to stay' and \zh{学} 'to study', both inherited from Middle Chinese, \emph{ljuw} and \emph{haewk} (Guangyun, in Baxter's transcription), which glosses the first 'to stay' and 'to stop'.
\end{cousinweb}

\begin{grammarlens}[title={What you've built}]
One more word for talking about the past.
\end{grammarlens}

\subsection*{Your turn}
\begin{itemize}
\raggedright
  \item \emph{Say it:} \emph{liúxué}
  \item \emph{Say it:} \emph{liúxué}, once more
  \item \emph{Say it:} \emph{wǒ xiǎng qù Zhōngguó liúxué} — I'd like to go to China to study.
  \item \emph{Recall:} say the Mandarin for know-how, then the Mandarin for at one time, then say \emph{liúxué} again
\end{itemize}

\begin{tcolorbox}[breakable,colback=teal!4,colframe=teal!35!black,title={Before you move on}]
What does liúxué mean? (\textquotedblleft{}To study abroad\textquotedblright{}.) Say it once more.
\end{tcolorbox}
\section[\texorpdfstring{\zh{国外} (guówài)}{国外 (guówài)}]{\zh{国外} (guówài) — abroad; overseas}

\label{lesson:ZH-C261-guowai}



\begin{quote}
Before the new one: say the Mandarin for at one time, then the Mandarin for to study abroad.
\end{quote}

\subsection*{You'll want to know: \zh{国外}}
\textbf{\zh{国外}} — \emph{guówài} — \textquotedblleft{}abroad; overseas\textquotedblright{}.

\begin{quote}
\textbf{\zh{你去过国外吗}} \emph{nǐ qù guo guówài ma} — Have you ever been abroad?
\end{quote}

\begin{cousinweb}
\textbf{\zh{国外}} (\emph{guówài}) is built from \zh{国} 'a country' and \zh{外} 'outside', both inherited from Middle Chinese, \emph{kwok} and \emph{ngwajH} (Guangyun, in Baxter's transcription), which glosses the second 'outside' and 'far'.
\end{cousinweb}

\begin{grammarlens}[title={What you've built}]
One more word for talking about the past.
\end{grammarlens}

\subsection*{Your turn}
\begin{itemize}
\raggedright
  \item \emph{Say it:} \emph{guówài}
  \item \emph{Say it:} \emph{guówài}, once more
  \item \emph{Say it:} \emph{nǐ qù guo guówài ma} — Have you ever been abroad?
  \item \emph{Recall:} say the Mandarin for at one time, then the Mandarin for to study abroad, then say \emph{guówài} again
\end{itemize}

\begin{tcolorbox}[breakable,colback=teal!4,colframe=teal!35!black,title={Before you move on}]
What does \zh{国外} mean? (\textquotedblleft{}Abroad; overseas\textquotedblright{}.) Say it once more.
\end{tcolorbox}
\section[tǐyàn]{tǐyàn — to experience for yourself}

\label{lesson:ZH-C261-tiyan}



\begin{quote}
Before the new one: say the Mandarin for to study abroad, then the Mandarin for overseas.
\end{quote}

\subsection*{You'll want to know: tǐyàn}
\emph{tǐyàn} — \textquotedblleft{}to experience for yourself\textquotedblright{}.

\begin{quote}
\emph{tāmen qù nóngcūn tǐyàn shēnghuó} — They're going to the countryside to experience life there.
\end{quote}

Its characters come later; for now it is a word for the ear.

\begin{cousinweb}
\emph{tǐyàn} is built from \zh{体} 'the body' and \emph{yàn} 'to verify', both inherited from Middle Chinese, \emph{thejX} and \emph{ngjemH} (Guangyun, in Baxter's transcription), which glosses the first 'the body' and the second 'to verify'.
\end{cousinweb}

\begin{grammarlens}[title={What you've built}]
One more word for talking about the past.
\end{grammarlens}

\subsection*{Your turn}
\begin{itemize}
\raggedright
  \item \emph{Say it:} \emph{tǐyàn}
  \item \emph{Say it:} \emph{tǐyàn}, once more
  \item \emph{Say it:} \emph{tāmen qù nóngcūn tǐyàn shēnghuó} — They're going to the countryside to experience life there.
  \item \emph{Recall:} say the Mandarin for to study abroad, then the Mandarin for overseas, then say \emph{tǐyàn} again
\end{itemize}

\begin{tcolorbox}[breakable,colback=teal!4,colframe=teal!35!black,title={Before you move on}]
What does tǐyàn mean? (\textquotedblleft{}To experience for yourself\textquotedblright{}.) Say it once more.
\end{tcolorbox}
